\section{Invariant Nash Equilibrium and Tight Constants}

Let $X=\prod_i\Delta(A_i)$ be the mixed-profile space and
$X^G=\{\sigma:g\sigma=\sigma\ \forall g\in G\}$.  The fixed set is nonempty, compact, and convex: averaging any profile over $G$ produces an element of $X^G$.

\begin{theorem}[Invariant Nash equilibrium]
\label{thm:supp-invariant-ne}
Every finite exactly $G$-invariant game has a Nash equilibrium in $X^G$.
\end{theorem}
\begin{proof}
For precision $\lambda>0$, define the logit response
\[
F_i^\lambda(\sigma)(a_i)
=
\frac{\exp(\lambda u_i(a_i,\sigma_{-i}))}
{\sum_{b_i}\exp(\lambda u_i(b_i,\sigma_{-i}))}.
\]
Exact invariance implies equivariance:
$F^\lambda(g\sigma)=gF^\lambda(\sigma)$.  Thus $F^\lambda$ maps $X^G$ into itself.  Brouwer gives a fixed point $\sigma^\lambda\in X^G$.  By compactness, along a sequence $\lambda_k\to\infty$ we have
$\sigma^{\lambda_k}\to\sigma^*\in X^G$.

Suppose an action $a_i$ in the support of $\sigma_i^*$ is worse than an action $b_i$ against $\sigma_{-i}^*$ by margin $\gamma>0$.  By continuity, the margin is at least $\gamma/2$ along the sequence for all large $k$.  At a logit fixed point,
\[
\frac{\sigma_i^{\lambda_k}(a_i)}
{\sigma_i^{\lambda_k}(b_i)}
\le e^{-\lambda_k\gamma/2}\to0,
\]
which contradicts positive limiting mass on $a_i$.  Every supported action is a best response, so $\sigma^*$ is Nash.
\end{proof}

\begin{theorem}[Nash transfer]
\label{thm:supp-ne-transfer}
If $\devdist(\Gamma,\widehat\Gamma)\le\delta$ and $\sigma$ is an $\epsilon$-Nash equilibrium of $\widehat\Gamma$, then $\sigma$ is an $(\epsilon+\delta)$-Nash equilibrium of $\Gamma$.
\end{theorem}
\begin{proof}
For every player and deviation $\tau_i$, Lemma~\ref{lem:supp-mixed} gives
\[
u_i(\tau_i,\sigma_{-i})-u_i(\sigma)
\le
\widehat u_i(\tau_i,\sigma_{-i})-
\widehat u_i(\sigma)+\delta
\le\epsilon+\delta.
\]
\end{proof}

Combining Theorems~\ref{thm:supp-invariant-ne} and
\ref{thm:supp-ne-transfer} with a nearest symmetric surrogate proves that every finite game has a $G$-respecting $\delta_G$-Nash equilibrium.

\subsection{Tightness of the Transfer Coefficient}

Fix $k\ge2$.  Player 1 has $k$ actions, obtains payoff one from action 1 and zero from every other action, independently of player 2; player 2 is strategically irrelevant.  Let $G_k$ cycle player 1's actions.  Every invariant surrogate equalizes those $k$ action payoffs.  Since the original incentive difference between action 1 and another action is one, every invariant surrogate has strategic distance at least one, while the zero surrogate attains one.  Hence $\delta_{G_k}=1$.

The only invariant strategy of player 1 is uniform.  Its payoff is $1/k$, while deviating to action 1 gives one, so its regret is $1-1/k$.  Therefore no universal equilibrium-transfer theorem can replace the additive coefficient one by a smaller constant.

\section{Group Averaging}

Define the Reynolds projection and direct orbit discrepancy by
\[
\mathcal P_G\Gamma=|G|^{-1}\sum_{g\in G}g\Gamma,
\qquad
\eta_G(\Gamma)=\max_{g\in G}\devdist(\Gamma,g\Gamma).
\]

\begin{theorem}[Factor-two averaging]
\label{thm:supp-reynolds}
For every finite game and finite relabeling group,
\[
\delta_G(\Gamma)
\le
\devdist(\Gamma,\mathcal P_G\Gamma)
\le
\eta_G(\Gamma)
\le
2\delta_G(\Gamma).
\]
\end{theorem}
\begin{proof}
The average is invariant, proving the first inequality.  Convexity of the seminorm gives
\[
\devdist(\Gamma,\mathcal P_G\Gamma)
\le
|G|^{-1}\sum_g\devdist(\Gamma,g\Gamma)
\le\eta_G(\Gamma).
\]
For any invariant surrogate $\widehat\Gamma$ and $g\in G$,
\begin{align*}
\devdist(\Gamma,g\Gamma)
&\le
\devdist(\Gamma,\widehat\Gamma)
+
\devdist(\widehat\Gamma,g\Gamma)\\
&=
\devdist(\Gamma,\widehat\Gamma)
+
\devdist(g\widehat\Gamma,g\Gamma)\\
&=2\devdist(\Gamma,\widehat\Gamma).
\end{align*}
Minimize over invariant surrogates and maximize over $g$.
\end{proof}

The average is computable without enumerating $G$: each payoff-coordinate orbit receives its arithmetic mean.  If the original game lies in a relabeling-invariant convex linear class, such as zero-sum or common-payoff games, the average remains in that class.

\subsection{Zero-Sum Sharpness Family}

Fix $d\ge2$.  The row player has actions $R=\{0,\ldots,d-1\}$ and the column player has actions
$C=\{0,1\}\times\{0,\ldots,d-1\}$.  Let $G_d$ cycle all row actions and independently cycle column actions within each block.  Define the row payoff, independent of the within-block column index, by
\[
A_d(s,(c,t))
=
\begin{cases}
c,&s=0,\\
1-c,&s\ne0.
\end{cases}
\]
The column payoff is $-A_d$.

The zero game is an invariant zero-sum surrogate at strategic distance one.  Conversely, invariance equalizes the row payoffs of the distinguished row and any other row at a fixed column, while their original difference is one.  Thus
\[
\delta_{G_d}(\Gamma_d)
=
\delta_{G_d,\mathcal Z}(\Gamma_d)
=1.
\]
Under group averaging, row-payoff means in the two column blocks are $(d-1)/d$ and $1/d$.  The distinguished row therefore has residual levels $-(d-1)/d$ and $(d-1)/d$, whose range is
\[
2-\frac{2}{d}.
\]
All other residual slices are no larger.  Hence the factor-two approximation is asymptotically tight even when both the original game and the nearest structured surrogate are zero-sum.
